\par\HistoricalFoot{}{\(7^*\)}
\indent Unter allen die Grundlagen betreffenden Fragen erscheint sowohl in ihrer
Bedeutung für \textsc{Gauss’} innere Entwicklung, als auch in bezug auf ihre allgemeine
Tragweite, die Lehre von den \textls[120]{komplexen Zahlen} als die wichtigste.
— Die Zweckmässigkeit, ja die Unentbehrlichkeit der damals noch sogenannten
imaginären Zahlen stand für \textsc{Gauss} wohl schon fest, seit er die Theorie der
Kreisteilung entwickelt hatte, was nach einer Äusserung von 1819 (Werke
X, 1, S. 125) bis in sein erstes Göttinger Studiensemester zurückreicht. —
Auch die älteste Stelle, der Art. 91 der \textit{Disquisitiones arithm.}, wo \textsc{Gauss} eine
allgemeine Theorie der imaginären Grössen in Aussicht stellt, stammt (vergl.
\textsc{Bachmann}, Abh. 1 dieses Bandes, Nr. 2.) aus dem Jahre 1796; er sagt daselbst
aus Anlass der imaginären Indizes für Zahlen der Form \(8n+3,\ 8n+7\),
dass er sich deren Erörterung für eine andere Gelegenheit vorbehält, wo er
»vielleicht ausführlich die Theorie der imaginären Grössen zu behandeln
unternehmen würde, die nach seinem Urteil bisher von niemand auf klare
Begriffe zurückgeführt worden sei«. Und dieses Desiderium einer ausreichenden
Rechtfertigung für die Zulässigkeit der imaginären Grössen als »gleichberechtigter
Bürger im Reiche der Zahlen« hat \textsc{Gauss} immer wieder betont, auch
hat der Mangel einer solchen Rechtfertigung ihn veranlasst, noch in der Dissertation
nicht nur den Gebrauch jener Grössen zu vermeiden, sondern auch mit
skeptisch kritischen Bemerkungen gegen ihren rein formalen Gebrauch nicht
zurückzuhalten. — Für seine Einstellung zu den komplexen Grössen dürfte
aber die Entdeckung der doppelten Periodizität der lemniskatischen Funktionen
von ausschlaggebendem Einfluss gewesen sein (Mai 1797, vergl. oben S. 33 ff.),
die ihm die Bedeutung der komplexen Grössen in zwiefacher Hinsicht enthüllt
hat, nämlich als \textls[120]{Zahlen} in der Arithmetik und als \textls[120]{Variable} in der
Analysis\footnote{Vergl. \textsc{Bachmann}, Abh. 1 dieses Bandes, S. 55.}.

Die Erscheinung, dass der sinus lemniscaticus seinen Wert nicht verändert,
wenn das Argument um \((m+n\sqrt{-1})\,2\varpi\) vermehrt wird, wo \(m,n\) reelle
ganze Zahlen bedeuten, dass also für diese Funktion die komplexen ganzen
Zahlen \(m+n\sqrt{-1}\) dieselbe Rolle spielen, wie für den gewöhnlichen Sinus die\par
\clearpage\SourceLeaf{137}\ArithmeticHead{52}{L}{L. SCHLESINGER, ÜBER GAUSS’ FUNKTIONENTHEORETISCHE ARBEITEN.}
\noindent reellen, muss auf \textsc{Gauss} einen tiefen Eindruck gemacht haben, zumal auch
bei den Problemen der Multiplikation und der Teilung sich das analoge Entsprechen
zur Geltung bringen liess. Wir werden also \textsc{Jacobi} beistimmen
müssen, wenn er\footnote{\textsc{C. G. J. Jacobi}, \textsc{Crelles} Journal 19, \textsc{Jacobis} Werke VI, S. 275 ff.} über die von \textsc{Gauss} bewirkte Einführung der komplexen
ganzen Zahlen in die Theorie der biquadratischen Reste sich wie folgt ausspricht:

»Aber wie einfach jetzt auch eine solche Einführung der komplexen Zahlen
als \textls[120]{Moduln} erscheinen mag, so gehört sie nichtsdestoweniger zu den tiefsten
Gedanken der Wissenschaft; ja ich glaube nicht, dass zu einem so verborgenen
Gedanken die Arithmetik allein geführt hat, sondern, dass er aus dem Studium
der elliptischen Transzendenten geschöpft ist, und zwar der besonderen Gattung
derselben, welche die Rektifikation vom Bogen der Lemniscata gibt. In der
Theorie der Vervielfachung und Teilung von Bogen der Lemniscata spielen
nämlich die komplexen Zahlen von der Form \(a+b\sqrt{-1}\) genau die Rolle gewöhnlicher
Zahlen. Wie man durch rationale Ausdrücke die trigonometrischen
Funktionen des \(n\)-fachen Kreisbogens darstellt, so kann man vermittelst rationaler
Formeln den Bogen der Lemniscata mit einer komplexen Zahl \(a+b\sqrt{-1}\)
multiplizieren; wie man den Kreisbogen durch Auflösung einer Gleichung
vom \(n\)-ten Grade in \(n\) Teile teilt, so teilt man den Bogen der Lemniscata in
\(a+b\sqrt{-1}\) Teile durch Auflösung einer Gleichung vom Grade \(aa+bb\). .\ .\ .
So wird man bei Untersuchung jener besonderen Gattung elliptischer Integrale,
wenn man nur einigermassen in ihre Natur eindringt, mit Notwendigkeit
darauf hingedrängt, die Zahlen \(a+b\sqrt{-1}\) als \textls[120]{Divisoren} einzuführen.
Mögen nun auch jene Untersuchungen der Integralrechnung viel komplizierter
und schwieriger erscheinen, als jener einfache Gedanke der Zahlenlehre, so
ist es doch nicht immer das Einfache, welches sich zuerst darbietet. \textsc{Gauss}
versichert in den \textit{Disquisitiones arithmeticae} (Sect. VII, Werke I, S. 413), die
Methode seiner Kreisteilung auf die Teilung der ganzen Lemniscata anwenden
zu können, und verspricht hierüber ein \textls[120]{amplum opus} zu einer Zeit, in
welcher er sich sicher noch nicht, seinen eigenen späteren Angaben zufolge,
mit den biquadratischen Resten beschäftigt hatte. Auch ist es nicht unwahrscheinlich,
dass er die Fundamentaltheoreme über biquadratische Reste aus
dieser Quelle geschöpft hat.«\par
\clearpage\SourceLeaf{138}\ArithmeticHead{53}{R}{III d. GRUNDLAGEN DER ANALYSIS. KOMPLEXE GRÖSSEN.}
Noch unmittelbarer ergab sich, wie wir gesehen haben, aus der Erkenntnis
der doppelten Periodizität der lemniskatischen Funktionen die Notwendigkeit,
das Gebiet einer \textls[120]{veränderlichen} Grösse dadurch zu erweitern, dass dieser
Grösse beliebige komplexe Werte beigelegt werden.

Schon \textsc{Euler} hatte zwar die Einführung der komplexen Variabeln in
formaler Weise vollzogen, indem er einerseits den Zusammenhang zwischen
den Kreisfunktionen und der Exponentialfunktion begründet, andererseits die
Logarithmen negativer und imaginärer Zahlen untersucht hat\footnote{Über \textsc{Eulers} Untersuchungen zur Theorie der Funktionen einer komplexen Variabeln berichtet
\textsc{Stäckel}, Bibliotheca Mathem. (III) 1, 1900, S. 113 ff.}. Für \textsc{Gauss}
handelte es sich aber darum, durch Eingehen auf die »wahre Metaphysik der
imaginären Grössen« ihre reale Existenz und damit die Berechtigung, sich
ihrer in der Analysis zu bedienen, definitiv zu begründen. Vor der Öffentlichkeit
geschieht dies erst im Jahre 1831\footnote{\textit{Theoria residuorum biquadraticorum, Commentatio secunda}, Artikel 30—38, Werke II, S. 102—110 und \textit{Anzeige} dieser Abh. ebenda S. 169 ff.; vergl. besonders S. 175 ff.}, weil — wie \textsc{Gauss} sich dort
ausdrückt — »es bisher an einer Veranlassung gefehlt hat«, er fügt aber gleich
hinzu, dass »aufmerksame Leser die Spuren davon in der 1799 erschienenen
Schrift über die Gleichungen und in der Preisschrift über die Umbildung der
Flächen leicht wiederfinden werden.«

Die erstere Schrift fällt, wie wir wissen, ihrer Entstehungszeit nach in
das Jahr 1797, wir werden also für dieses Jahr einerseits die Entdeckung der
Darstellung komplexer Grössen durch die Punkte einer Ebene, andererseits,
als erste Anwendung dieser Darstellung, den geometrischen Beweis der Wurzelexistenz
anzusetzen haben.

Was zunächst den ersten Punkt betrifft, so heisst es in der \textit{Anzeige}
(Werke II, S. 175): »Die Realität der negativen Zahlen ist hinreichend gerechtfertigt,
da sie in unzähligen Fällen ein adäquates Substrat findet, .\ .\ . allein
die den reellen Grössen gegenübergestellten imaginären .\ .\ .\ . sind noch immer
weniger eingebürgert als nur geduldet. .\ .\ .\ . Der Verfasser hat diesen hochwichtigen
Teil der Mathematik seit vielen Jahren aus einem verschiedenen
Gesichtspunkte betrachtet, wobei den imaginären Grössen ebensogut ein Gegenstand
unterlegt werden kann, wie den negativen.« Und in der Abhandlung
selbst liest man, nachdem im Art. 30 die komplexen ganzen Zahlen einge-\par
\clearpage\SourceLeaf{139}\ArithmeticHead{54}{L}{L. SCHLESINGER, ÜBER GAUSS’ FUNKTIONENTHEORETISCHE ARBEITEN.}
\noindent führt worden sind\footnote{Vergl. übrigens auch das aus 1811 stammende Nachlassstück, Werke II, S. 313.}, im Art. 38: »In den Grenzen dieser Untersuchung wird
es angemessen sein, anzugeben, wie man das Gebiet der komplexen Grössen
der Anschauung unterwerfen kann. So wie man eine jede reelle Grösse durch
eine von einem Anfangspunkte einer beiderseits unendlichen Geraden aus genommene
Strecke, die durch eine beliebige Einheitsstrecke gemessen wird,
oder auch durch deren andern Endpunkt darstellen kann, so dass die Punkte
auf der einen Seite vom Anfangspunkte aus die positiven, die auf der andern
Seite, die negativen Grössen repräsentieren, so kann jede komplexe Grösse
\(x+yi\) durch den Punkt einer unendlichen Ebene dargestellt werden, dessen
Abszisse \(=x\), dessen Ordinate (zu einer Seite der Abszissenlinie positiv, zur
andern negativ genommen) \(=y\) ist. Dies festgesetzt, kann man sagen, dass
eine beliebige komplexe Grösse die Ungleichheit zwischen der Lage des Punktes,
auf den sie sich bezieht, und der des Anfangspunktes misst, wenn man als
positive Einheit \(+1\) eine beliebige bestimmte Abweichung nach einer beliebigen
bestimmten Richtung bezeichnet, als negative Einheit \(-1\) die gleiche
Abweichung nach der entgegengesetzten Richtung und als imaginäre Einheiten
\(+i,-i\) die gleich grossen Abweichungen nach zwei seitlich senkrechten Richtungen.
So wird die Metaphysik der Grössen, die wir imaginär nennen, vortrefflich
illustriert. Sind \(m,m'\) zwei komplexe Grössen, \(M,M'\) die Punkte,
deren relative Lage zum Anfangspunkte (0) sie ausdrücken, so wird \(m-m'\)
nichts anderes sein, als die relative Lage von \(M\) zu \(M'\), während \(mm'\) die
Lage eines Punktes \(N\) relativ zu (0) darstellt, die durch die Lage von \(M\) relativ
zu (0) ebenso bestimmt wird, wie die Lage von \(M'\) durch die Lage des
Punktes, dem \(+1\) entspricht, so dass man nicht unangemessen sagen kann,
dass die Lagen der Punkte, die den komplexen Grössen \(mm',m,m',1\) entsprechen,
eine \textls[120]{Proportion} bilden. Eine eingehendere Behandlung dieser
Sache behalten wir uns für eine andere Gelegenheit vor.« Hält man hierzu
noch die Stelle der \textit{Anzeige} (Werke II, S. 177), wo gesagt wird, dass die bisher
unbestimmt gebliebenen »imaginären Einheiten« \(+i,-i\) in dem angegebenen
Sinne mittlere Proportionalgrössen zwischen \(+1\) und \(-1\) sind, also dem Zeichen
\(\sqrt{-1}\) entsprechen, so kann man sagen, dass \textsc{Gauss} in diesen wenigen Zeilen
in knapper, aber alles wesentliche erschöpfender Weise die Darstellung der
komplexen Grössen als Vektoren in der Ebene und die Gesetze des Rechnens\par
\clearpage\SourceLeaf{140}\ArithmeticHead{55}{R}{III d. GRUNDLAGEN DER ANALYSIS. KOMPLEXE GRÖSSEN.}
\noindent mit solchen Vektoren gegeben hat. Dabei legt \textsc{Gauss}, wie aus den weiteren
Ausführungen der \textit{Anzeige} und besonders scharf aus der, aus dem Jahre 1825\footnote{Vergl. den aus etwa derselben Zeit stammenden Brief an \textsc{Hansen}, Werke XII, S. 6—9.}
stammenden, nachgelassenen Notiz (Werke X, 1, S. 396, 397) hervorgeht, den
Hauptnachdruck auf den Umstand, dass, während bei den reellen Zahlen das
zu ihrer Versinnlichung dienende Substrat aus Dingen besteht, wo jedes, wie
die Punkte einer Linie »nur zu zweien ein Verhältnis der Ungleichheit hat«,
bei den komplexen Zahlen Dinge in Betracht kommen, die, wie die Punkte
einer Fläche, »zu dreien in einem Wechselverhältnis stehen, und wo es zwischen
den Verhältnissen ein Verhältnis gibt«. Und nun folgt die bereits oben angeführte
Stelle von der Wichtigkeit, »die Theorie des Gegensatzes zur Klarheit
zu bringen ohne Grösse« also, wie wir sagen würden, ohne Metrik, nur durch
die ein bezw. zweidimensionale Anordnung, ein Prinzip, dem die oben geschilderte
Einführung des Grenzbegriffs durchaus angepasst erscheint. — Es
ist auch keine Spur dafür vorhanden, dass \textsc{Gauss} jemals für den Grenzbegriff
im komplexen Gebiete die uns seit \textsc{Cauchy} so geläufige Definition mit Hilfe
des Abstandes (absoluten Betrages) benutzt haben würde, vielmehr hat er eben
die Existenz des Grenzwerts einer komplexen Zahlenfolge einfach in der Existenz
der Grenzwerte für die Folgen der beiden Koordinaten gesehen.

Erst nach der Erörterung dieser abstrakten Relationen folgt in der \textit{Anzeige}
(Werke II, S. 176, 177) die anschauliche »Darstellung im Raume«, wo
dann »kein Grund vorhanden ist, die Symbole der Gegenstände anders als
quadratisch anzuordnen« und man auch ganz »nach Gefallen« den sich auf
\(+i\) beziehenden Punkt »rechts oder links« von der die Punkte \(+1\) und \(-1\)
verbindenden Linie nehmen kann\footnote{Auf die Erörterung über die Unterscheidung zwischen rechts und links, die in der \textit{Anzeige} folgt,
und die damit zusammenhängende Auseinandersetzung mit \textsc{Kant} kommen wir an späterer Stelle (Abschnitt
VIII, c)) zurück.}.
